% SOURCE_AUTHORITY: sga5_fr_workpass.tex, lines 8965--8978 (LF-normalized unit).
% SOURCE_SHA256: 791F4EFFC5E02832D5D77ED03518C8156D6F07E4C8238B03545DB93D883FBB28
\subsection*{3. Clase de cohomología \(\ell\)-ádica asociada a un ciclo}

\subsubsection*{3.1}
Si \(X\) es un esquema propio, o liso y de tipo finito, sobre un cuerpo algebraicamente cerrado \(k\), se sabe que, para todo haz constructible \(M\), abeliano y de torsión prima con la característica de \(k\), los grupos de cohomología
\[
H^i(X,M)
\]
son grupos abelianos de longitud finita. Esto permite (2.2.1, Observación) definir, para todo \(\mathbb Z_\ell\)-haz \(F\) (\(\ell\) primo con la característica de \(k\)), un sistema proyectivo \((AR,\ell)\)-ádico noetheriano de grupos abelianos
\[
H^i(X,F)
\]
con las propiedades habituales. En particular, el método utilizado en el número 2 permite definir el producto cup, el homomorfismo de imagen inversa para un morfismo de esquemas sobre \(k\) y el homomorfismo de imagen directa para un morfismo propio de esquemas lisos sobre \(k\) (cf. Exp. IV para las definiciones en el caso de los haces de torsión).

Por otra parte, se sabe (1.2.4) que la categoría de los \(\mathbb Z_\ell\)-haces sobre un cuerpo algebraicamente cerrado es equivalente a la categoría de los \(\mathbb Z_\ell\)-módulos de tipo finito, y que la equivalencia viene dada por el funtor límite proyectivo. En la práctica será conveniente identificar el \(\mathbb Z_\ell\)-módulo de tipo finito \(H^i(X,F)\) con el \(\mathbb Z_\ell\)-haz correspondiente.
